% Generated draft: requires mathematical and source review.
\documentclass[11pt]{article}
\usepackage[margin=25mm]{geometry}
\usepackage{fontspec}
\setmainfont{Latin Modern Roman}
\usepackage{amsmath,amssymb,mathtools,unicode-math}
\setmathfont{Latin Modern Math}
\usepackage{enumitem,needspace,xurl,hyperref}
\setlength{\parindent}{0pt}
\setlength{\parskip}{5pt}
\setlength{\emergencystretch}{3em}
\newcommand{\N}{\mathbb{N}}
\newcommand{\Z}{\mathbb{Z}}
\newcommand{\Q}{\mathbb{Q}}
\newcommand{\R}{\mathbb{R}}
\newcommand{\C}{\mathbb{C}}
\newcommand{\F}{\mathbb{F}}
\newcommand{\A}{\mathbb{A}}
\newcommand{\PP}{\mathbb{P}}
\newcommand{\E}{\mathbb{E}}
\newcommand{\eps}{\varepsilon}
\newcommand{\dd}{\,\mathrm{d}}
\begin{document}
\section*{UnsolvedMath ID: 701497}
\textbf{Source problem code:} catalogue entry
\subsection{Definitions and mathematical statement}
Let \(G\) be a finite group. A subgroup \(M<G\) is maximal if there is no subgroup strictly between \(M\) and \(G\). A complement to \(M\) in \(G\) is a subgroup \(H\leq G\) such that
\[
G=MH,\qquad M\cap H=1.
\]
Neither subgroup is required to be normal. Say that \(G\) has complemented maximal subgroups if
\[
\forall M<G\text{ maximal},\quad \exists H\leq G:\quad G=MH\ \text{ and }\ M\cap H=1.
\]
A group is soluble if its derived series, defined by \(G^{(0)}=G\) and \(G^{(i+1)}=[G^{(i)},G^{(i)}]\), reaches the trivial group. A composition series is a chain
\[
1=G_0\triangleleft G_1\triangleleft\cdots\triangleleft G_r=G
\]
whose factors \(G_i/G_{i-1}\) are nontrivial simple groups. Its nonabelian composition factors are those factors that are nonabelian, considered up to isomorphism.

\textbf{Problem (N. V. Maslova and D. O. Revin, 2014; Kourovka Notebook 18.68).} Determine all isomorphism types of nonabelian simple groups that occur as composition factors of finite nonsoluble groups with complemented maximal subgroups. Equivalently, determine the set
\[
\mathcal S=\left\{[S]_{\cong}:\begin{array}{l}
S\text{ is a finite nonabelian simple group, and there exists a finite}\\
\text{nonsoluble group }G\text{ with complemented maximal subgroups}\\
\text{such that }S\text{ is isomorphic to a composition factor of }G
\end{array}\right\}.
\]

\textbf{Supplied background and qualifications.} Write \(L_n(q)=\operatorname{PSL}_n(\F_q)\), where \(\F_q\) is the field with \(q\) elements. The source attributes to V. M. Levchuk and A. G. Likharev (2006) the assertion that the finite simple groups with complemented maximal subgroups are, up to isomorphism,
\[
L_2(7),\qquad L_2(11),\qquad L_5(2).
\]
A subgroup \(M\leq G\) is a Hall subgroup if \(\gcd(|M|,[G:M])=1\); having Hall maximal subgroups means that every maximal subgroup is a Hall subgroup. The source attributes to N. V. Maslova (2012) the assertion that the same three groups exhaust the finite simple groups with Hall maximal subgroups and the composition factors of groups with Hall maximal subgroups. It further attributes to N. V. Maslova and D. O. Revin (2013) the assertion that every maximal subgroup of a finite group with Hall maximal subgroups has a complement.

The supplied background does not explicitly restrict its references to simple groups and composition factors to the nonabelian ones. This qualification requires checking: abelian simple groups are cyclic of prime order, and such groups also have complemented maximal subgroups and Hall maximal subgroups. Accordingly, the background assertions are recorded as supplied attribution with unresolved scope; they are not imposed as additional hypotheses or asserted here as a classification of \(\mathcal S\).

The supplied registry labels the problem open. It also records an explicitly unverified AI solution claim that the factors have been described in a 2026 preprint attributed to R. Sater. Neither that claim nor the background literature has been independently verified for this draft.
\subsection{Short English statement}
Determine which finite nonabelian simple groups can occur as composition factors of a finite nonsoluble group in which every maximal subgroup has a complement. The supplied background classification statements have an unresolved qualification concerning abelian simple groups and factors. The registry labels the problem open and separately records an unverified AI solution claim.
\subsection{Sources}
\begin{itemize}
\item Local UnsolvedMath source registry, supplied complete record: permanent ID \(701497\), exact source problem code \texttt{catalogue entry}, title \emph{Kourovka Notebook Problem 18.68}. Proposed by N. V. Maslova and D. O. Revin in 2014.
\item \emph{The Kourovka Notebook (Unsolved Problems in Group Theory)}, problem 18.68, 18th issue (2014); supplied source also identifies the 21st issue edition (2026), \texttt{arXiv:1401.0300v48}. The registry describes the collection as Issues 1--21 (1965--2026), edited by E. I. Khukhro and V. D. Mazurov.
\item V. M. Levchuk and A. G. Likharev, \emph{Siberian Math. J.}, \textbf{47}, no. 4 (2006), 659--668. Supplied attribution for the simple-group classification concerning complemented maximal subgroups; article title and external link were not supplied.
\item N. V. Maslova, \emph{Siberian Math. J.}, \textbf{53}, no. 5 (2012), 853--861. Supplied attribution for the assertions concerning Hall maximal subgroups and composition factors; article title and external link were not supplied.
\item N. V. Maslova and D. O. Revin, \emph{Siberian Adv. Math.}, \textbf{23}, no. 3 (2013), 196--209. Supplied attribution for the implication from Hall maximal subgroups to complemented maximal subgroups; article title and external link were not supplied.
\item R. Sater, \emph{Preprint} (2026), \url{https://kourovkanotebookorg.wordpress.com/wp-content/uploads/2026/09/kourovka-18-68-v1.1.1.pdf}. The supplied background and status review explicitly identify this as an unverified AI solution claim, described as asserting that these factors have been determined. No verification or acceptance of that claim is asserted here.
\item Supplied status review dated \texttt{2026-10-16}: canonical status \texttt{open}, evidence category \texttt{unrefereed\_claim}. Its notes report a \(\#\) marker for an unverified AI solution claim in the October 2026 21st issue, \texttt{arXiv:1401.0300v48}. This is supplied dated curation, not an independently verified literature assessment.
\end{itemize}
\end{document}
